%!TEX root = ../atomic_jsc.tex
% LTeX: language=en
%
\section{Introduction}
%
Let $\K$ be a field and $\Indets$ be a countable (i.e., finite or countably infinite) set of variables.
Let $\monoid$ be a monoid acting injectively on $\Indets$
(i.e., the function $x\mapsto \gelem\cdot x$ is injective for every $\gelem\in\monoid$).
An ideal in $\poly{\K}{\Indets}$ is called $\monoid$-equivariant if it is invariant under the action of $\monoid$.
For instance, let $\inc{\N}$ be the monoid of strictly monotonic functions from $\N$ to itself and the action on $\Indets = \{x_1,x_2,\dots\}$ is defined as $\gelem \cdot x_i = x_{\gelem(i)}$ for $\gelem \in \inc{\N}$ and $x_i\in \Indets$.
Then the ideal $\idl[S]_0$ of polynomials in $\poly{\K}{\Indets}$ whose coefficients add up to zero is a $\inc{\N}$-equivariant ideal.
In \cite{COHEN67} Cohen showed that the classical Hilbert's basis theorem \arka{cite wiki} extends to $\inc{\N}$-equivariant ideals: every $\inc{\N}$-equivariant ideal in $\poly{\K}{\Indets}$ is generated by an orbit-finite set of polynomials.
For instance, $\idl[S]_0$ is generated by the set $\setof{x_i - 1}{i\in\N}$,
which forms a single orbit under the action of $\inc{\N}$.
Let $\symgr{\Indets}$ denote the group of all permutations of $\Indets$.
Since every $\symgr{\Indets}$-equivariant ideal is also $\inc{\N}$-equivariant,
as a corollary of Cohen's result one gets that $\symgr{\Indets}$-equivariant ideals are also finitely generated.
The results of Cohen were rediscovered in \cite{AH07} and generalised in \cite{HISU12,LauSno26,GHOLAS24}.

To explain these results we need to give a few definitions.
We say that the action $\monoid\actson\Indets$ is:
\begin{enumerate}
\item \intro{ordered} if it respects a total order $\prec$ on $\Indets$ (i.e., $\gelem \cdot x \prec \gelem \cdot y$ for every $x \prec y$ and $\gelem\in\monoid$),
\item \intro{well-ordered} if it respects a well-order on $\Indets$,
\item \intro{nice} if the the divisibility up to $\monoid$ relation ($\monelt[p] \gdivleq[\monoid] \monelt[q]$ if there exists $\pi\in\monoid$ such that $\pi\cdot\monelt[p]$ divides $\monelt[q]$) is a well-quasi-order on the set $\mon{\Indets}$ of monomials in $\Indets$,
\item \intro{nicely ordered} if it is \kl{nice} and \kl{ordered},
\item \intro{nicely well-ordered} if it is \kl{nice} and \kl{well-ordered}.
\end{enumerate}
%
Observe that an \kl{ordered} monoid action must be injective.
The action $\symgr{\Indets}\actson\Indets$ is \kl{nice} but not \kl{ordered}.
The action of $\inc{\N}$ mentioned above is \kl{nicely well-ordered}.
In \arka{refer sec} we discuss several examples of monoid actions and discuss which of the above properties they satisfy.


The article \cite{HISU12} showed proved $\monoid$-equivariant ideals are orbit-finitely generated when $\monoid\actson\Indets$ is \kl{nicely well-ordered},
and used it to prove the independent set conjecture in algebraic statistics.
This result was further generalised in \cite{GHOLAS24,LauSno26} which showed that:
\begin{enumerate}
\item $\monoid$-equivariant ideals are orbit-finitely generated when $\monoid\actson\Indets$ is \kl{nicely ordered} (i.e., they get rid of the well-foundedness condition on the ordering on variables).
\item Conversely, if $\monoid$-equivariant ideals are orbit-finitely generated then $\monoid\actson\Indets$ is \kl{nice}.
\end{enumerate}
%
It has been conjectured by Pouzet that if $\monoid\actson\Indets$ is \kl{nice},
then there exists another monoid $\mathcal{N}$ acting on $\Indets$ such that 
\begin{enumerate}
\item for every finite subset $S$ of $\Indets$ and $\gelem\in\mathcal{N}$,
there exists $\sigma\in\monoid$ such that $\sigma\cdot x = \gelem\cdot x$ for every $x\in S$.
\item $\mathcal{N}\actson\Indets$ is \kl{nicely ordered}. 
\end{enumerate}
%
Monoid actions $\monoid\actson\Indets$ satisfying these conditions will be called \intro{nicely orderable}.
%
\begin{remark}
It is tempting to conjecture that $\mathcal{N}$ can even be a sub-monoid of $\monoid$.
However this stronger conjecture would be false as there is no sub-monoid $\mathcal{N}$ of the group of finitary permutations of $\Indets$ (permutations $\pi$ that $\pi(x) \neq x$ only for finitely many $x\in\Indets$) that satisfy the two above conditions.
Observe however $\inc{\N}$ is a monoid that satisfies these conditions but is not a sub-monoid of the group of finitary permutations.
\end{remark}
%
\begin{example}
We give an example of a \kl{nicely ordered} action that is not well-ordered.
Let $\prec$ be a dense linear order on $\Indets$ without endpoints (such an order can be constructed from a bijection of $\Indets$ with the set of rational numbers), and $\monoid$ be the set of all strictly monotonic bijections from $\Indets$ to itself.
Clearly $\monoid\actson\Indets$ is not well-ordered.
However $\gdivleq[\monoid]$ is a well-quasi-order on $\mon{\Indets}$ due to Higman's lemma.
\arka{we discuss a proof scheme to prove niceness later..}
\end{example}
%
The result of \cite{HISU12} was extended in \cite{BRDR11} which defined Gr\"{o}bner bases for $\monoid$-equivariant ideals and showed that they are orbit-finite and computable by an appropriate extension of the classical Buchberger's 
\arka{gave 'a' definition. Mention nicely well-ordered case before. List the contributions}
algorithm in the equivariant setting when $\monoid\actson\Indets$ is \kl{nicely well-ordered} and satisfies certain computability conditions that are necessary for running the equivariant Buchberger's algorithm.
Our paper extends the results of \cite{GHOLAS24,LauSno26} and generalise the results of \cite{BRDR11} to the case when the action of $\monoid\actson\Indets$ is \kl{nicely ordered}.
%
\subsection{Contributions}
\arka{cite theorems and sections. Also list the contributions}

As our first main result we define Gr\"{o}bner bases for $\monoid$-equivariant ideals and show that when $\monoid\actson\Indets$ is \kl{nicely ordered} and satisfies certain computability criterions,
every $\monoid$-equivariant ideal has an orbit-finite Gr\"{o}bner basis and
there is an algorithm that takes an orbit-finite set $H\subseteq\poly{\K}{\Indets}$ as input and outputs an orbit-finite Gr\"{o}bner base for $\monoid$-equivariant ideal generated by $H$ (an orbit-finite set
$(\monoid\cdot f_1)\cup\dots\cup(\monoid\cdot f_n)$ is represented by the set $\set{f_1,\dots,f_n}$).

As a corollary of our results we get that the $\monoid$-equivariant ideal membership problem is decidable.
Moreover, given orbit-finite Gr\"{o}bner bases for $\monoid$-equivariant ideals $\idl[I]$ and $\idl[J]$ one can compute orbit-finite Gr\"{o}bner bases for $\monoid$-equivariant ideals generated by $\idl[I]\cup\idl[J]$, $\idl[I]\cap\idl[J]$ and $\idl[I]\idl[J]$.

We complement our positive results by giving a sufficient condition on $\monoid\actson\Indets$ for the undecidability of the $\monoid$-equivariant ideal membership problem.
This sufficient condition is satisfied by most common actions $\monoid\actson\Indets$ which are not \kl{nice} (and hence there are $\monoid$-equivariant ideals that are not finitely generated). 

Finally, we sketch some applications of our results for solving the membership problem for orbit-finitely spanned vector spaces that has received the attention of the theoretical computer science community in recent years.
\arka{cite papers}
%
\subsection{Comparison with previous work}
%
We quickly describe how we improve the state of the art on Gr\"{o}bner bases by comparing our results to our immediate predecessors \cite{GHOLAS24,BRDR11}.

In addition to proving orbit-finite generation of $\monoid$-equivariant ideals \cite{GHOLAS24} also show that for some specific \kl{nicely ordered} actions $\monoid\actson\Indets$,
one can reduce the $\monoid$-equivariant ideal membership problem to the $\mathcal{N}$-equivariant ideal membership problem for some nicely well-ordered action $\mathcal{N}\actson\Indets$
(we point the reader to \arka{refer} for more details).
However, it is unclear whether their technique can be generalised to the class of all \kl{nicely ordered} actions,
at least the authors of \cite{GHOLAS24} do not claim that to be the case.
In particular, it seems difficult to extend the technique of \cite{GHOLAS24} to actions $\monoid\actson\Indets$ where $\Indets$ is the structure of a densely coloured dense linear order or the ordered dense meet tree and $\monoid$ is the group of automorphisms of the structure.\arka{refer to def}
Both of these actions are \kl{nicely ordered}.
 
Furthermore, the reduction to the nicely well-ordered actions only solves the ideal membership problem.
Extending the notion of Gr\"{o}bner base to the general setting of \kl{nicely ordered} actions allows us to also show effectivity of simple operations as sum and intersection of equivariant ideals,
which would not be possible using the results of \cite{GHOLAS24}.
At the same time it opens up the opportunity to extend complicated operations such as computing the primary decomposition and radical of ideals to the equivariant setting. 

Let us now explain the difficulty in generalising the concepts of \kl{Gr\"{o}bner bases} and \kl{Buchberger's algorithm} from \kl{nicely well-ordered} actions to \kl{nicely ordered actions}.
When the action $\monoid\actson\Indets$ preserves an well-order $\prec$ on $\Indets$,
one can define a well-order on $\mon{\Indets}$ which is compatible with the monoid action and divisibility,
by extending $\prec$ to $\mon{\Indets}$ with some suitable ordering such as co-lex or grevlex \arka{refer to defintions}.
However the same is not possible for all \kl{nicely ordered} actions.
%In the case when the action $\monoid\actson\Indets$ is \kl{nicely ordered} and neither the order $\prec$ preserved by $\monoid\actson\Indets$ nor its reverse is well-founded,
%it is not possible to define a well-order on $\Indets$ which is compatible with the action as any such order must agree with $\prec$ or its reverse
%(for instance co-lex agrees with $\prec$ and lex agrees with its reverse
%\arka{this is not always true. But it is true for ordered atoms. Change later as well}).
This means that sequence of reductions that merely decreases a ordering of monomials based on $\prec$ may not terminate.
%
\begin{example}
For example, consider the action $\inc{\OrderAtoms}\actson\OrderAtoms$.
Let $\monlt$ be a monomial ordering on $\mon{\OrderAtoms}$ that is $\monoid$-equivariant.
Pick $x < y \in \OrderAtoms$.
If $x \monlt y$, then $x' < y'$ for every $x' < y'\in \OrderAtoms$, because for every such $x'$ and $y'$ there exists some $\pi\in\inc{\OrderAtoms}$ such that $\pi\cdot x = x'$ and $\pi\cdot y = y'$.
But then $\monlt$ agrees with $<$ and hence cannot be a well-order.
Similarly, if $x \monlt y$ then $\monlt$ agrees with the reverse of $\prec$ and hence cannot be a well-order.


Assume $(\Indets,\prec)$ is a dense linear order without endpoints.
If we try to reduce any polynomial $x\in\poly{\K}{\Indets}$ using the set of polynomials $\setof{y - z}{y \prec z}$ we get an infinite sequence
$x \xrightarrow{} x_1 \xrightarrow{} x_2 \xrightarrow{}\dots$
for some $x \succ x_1 \succ x_2 \succ\dots$. 
\end{example}
%
\arka{drop either equivariance or well-order.
We make an educated guess and choose to drop the latter.
But even then reductions are non-terminating}

To make reductions terminating we disallow reductions to introduce new variables.
But now the definition of Gr\"{o}bener basis has to modified accordingly:
for a set $\Basis$ of polynomials to be a Gr\"{o}bner basis of an ideal $\idl[I]$,
for every $f\in\idl[I]$ there must be an element $b\in\Basis$ which does not use any variable that is not used by $f$,
and the leading monomial of $b$ divides the leading monomial of $f$\arka{itemise}.
In this article we call them \kl{strong Gr\"{o}bner bases}\arka{is this a good name?}.
We use ideas from \cite{GHOLAS24} to show that when $\monoid\actson\Indets$ is \kl{nicely ordered},
every $\monoid$-equivariant ideal has an orbit-finite \kl{strong Gr\"{o}bner base}.
However, now the naive extension of the Buchberger's algorithm to the equivariant setting does not compute \kl{strong Gr\"{o}bner bases}.
We illustrate this through concrete examples in \arka{write and refer}.
As one of the main technical contribution of this paper,
we give a simple modification of the Buchberger's algorithm that allows us to compute \kl{strong Gr\"{o}bner bases}.
%
\subsection{Organisation of the article}
%
\subsection{Related research}
%
\subsubsection{Equivariant ideals}
%
Equivariant polynomial ideals have received significant attention from the mathematics community in recent years. 
Above we have mentioned the articles that study existence and computability of orbit-finite Gr\"{o}bner bases of equivariant ideals.
In addition to that,
there also has been a lot of research on the structure of equivariant ideals and varieties.
Hilbert series of equivariant ideals have been studied by \cite{NaRo17,KrLeSn17,MaNa21,Na21}
The articles \cite{DEKL15,NaSn20} have studied equivariant varieties.
Prime and radical equivariant ideals have been studied by \cite{NaSn21,KuRi25,Sn24}.
Finally \cite{SaSn16,SaSn19,Yu25,NSY25} have studied equivariant varieties.
%
\subsubsection{Orbit-finitely generated algebraic structures}
%
Our results also fit into the research agenda of studying algebraic structures built from orbit-finite sets.
These algebraic objects have recently caught the attention of the theoretic computer science community due to their connection with models of computations with infinite alphabets,
such as register automata and data Petri nets
\cite{Boj13,BoSt20,Ste24,BFKM24,YBK26,GHL22,GHL25,Gho25,ramsey}.
To properly discuss this connection we need to set up some background material which is done in \arka{refer}.
The algebraic structures and their connections to automata theory has been discussed in \arka{refer}.
%
\subsection{Finite groups???????}
%
%In this setting one considers algebraic structures (monoids, vector spaces, polynomial rings) built from $\omega$-categorical relational structures $\Indets$.
%
%For instance, if $\Indets$ is the structure of \kl{equality atoms} (infinite set without any additional structure) $\aut{\Indets} = \symgr{\Indets}$.
%And if $\Indets$ is the structure of \kl{ordered atoms} (dense linear order without endpoints), $\aut{\Indets}$ is the set of all order preserving bijections of $\Indets$.
%By Ryll-Nardzewski theorem,
%$\Indets$ being $\omega$-categorical is equivalent to the action $\group\actson\Indets$ being oligomorphic
%(i.e., the $\Indets^n$ has finitely many $\group$-orbits for every $n$).
%Now one studies algebraic structures built from the relational structure $\Indets$.
%Mainly three classes of algebraic structures have been studied in the literature,
%one of which are polynomial rings of the form $\poly{\K}{\Indets}$.
%The other two being
%\begin{enumerate}
%\item monoids of endomorphisms of $\Indets^n$ \cite{Boj13,BoSt20,Ste24},
%\item freely generated vector spaces $\K(\Indets^n)$ \cite{BFKM24,YBK26,GHL22,GHL25,Gho25,ramsey}.
%\end{enumerate}
%
%
\tableofcontents
%